\chapter{Ring of Integers of a Quadratic Field (Formalisation-Oriented)}

Let $d$ be a squarefree integer with $d \neq 1$.
We implement the ambient field as
\verb|QuadraticAlgebra| $\Q$ \verb|d 0|,
written in the project as \verb|QuadraticNumberFields.Qsqrtd|.

\begin{definition}\label{bp_qsqrtd}
  \lean{QuadraticNumberFields.Qsqrtd}
  \leanok
  The type \verb|Qsqrtd (d : \Q)| is the fixed Lean model of $\Q(\sqrt{d})$.
\end{definition}
\begin{proof}
  This is a type alias for \verb|QuadraticAlgebra| over $\Q$ with parameter $d$
  and zero linear term.
\end{proof}

\begin{lemma}\label{bp_trace}
  \lean{QuadraticNumberFields.Qsqrtd.trace_eq_two_re}
  \leanok
  For $x \in \Q(\sqrt{d})$, the trace is $2 \cdot \mathrm{re}(x)$.
\end{lemma}
\begin{proof}
  The explicit left-multiplication matrix in the basis $\{1, \sqrt{d}\}$ has
  diagonal entries both equal to the real part.
\end{proof}

\begin{definition}\label{bp_zsqrtd}
  \lean{QuadraticNumberFields.RingOfIntegers.Zsqrtd}
  \leanok
  The order $\Z[\sqrt{d}]$ is modelled by \verb|QuadraticAlgebra| $\Z$ \verb|d 0|.
\end{definition}
\begin{proof}
  This supplies an explicit ring type with the same quadratic relation as the
  field model.
\end{proof}

\begin{definition}\label{bp_zone}
  \lean{ZOnePlusSqrtOverTwo}
  \leanok
  When $d = 1 + 4k$, the order
  $\Z\left[\frac{1+\sqrt{d}}{2}\right]$ is modelled as
  \verb|QuadraticAlgebra| $\Z$ \verb|k 1|.
\end{definition}
\begin{proof}
  In this model the generator $\omega$ satisfies $\omega^2 = \omega + k$.
\end{proof}

\begin{lemma}\label{bp_trace_int}
  \lean{QuadraticNumberFields.RingOfIntegers.TraceNorm.exists_int_trace}
  \leanok
  \uses{bp_trace}
  If $x$ is integral over $\Z$, then its trace is an integer.
\end{lemma}
\begin{proof}
  This is the standard integrality property of trace, specialized to the
  quadratic-field model.
\end{proof}

\begin{lemma}\label{bp_norm_int}
  \lean{QuadraticNumberFields.RingOfIntegers.TraceNorm.exists_int_norm}
  \leanok
  If $x$ is integral over $\Z$, then its norm is an integer.
\end{lemma}
\begin{proof}
  This is the standard integrality property of norm.
\end{proof}

\begin{lemma}\label{bp_double_im}
  \lean{QuadraticNumberFields.RingOfIntegers.TraceNorm.exists_int_double_im}
  \leanok
  \uses{bp_trace_int, bp_norm_int}
  If $x$ is integral, then twice its imaginary coordinate is an integer.
\end{lemma}
\begin{proof}
  The proof combines integral trace, integral norm, and squarefree rigidity for
  the parameter $d$.
\end{proof}

\begin{lemma}\label{bp_halfint}
  \lean{QuadraticNumberFields.RingOfIntegers.exists_halfInt_with_div_four_of_isIntegral}
  \leanok
  \uses{bp_trace_int, bp_norm_int, bp_double_im}
  Any integral element of $\Q(\sqrt{d})$ has the form
  \[
  \frac{a' + b'\sqrt{d}}{2}
  \]
  with
  \[
  4 \mid (a'^2 - d b'^2).
  \]
\end{lemma}
\begin{proof}
  The integer witnesses for the trace, norm, and doubled imaginary part are
  assembled into the half-integer normal form.
\end{proof}

\begin{lemma}\label{bp_mod4_nonone}
  \lean{QuadraticNumberFields.RingOfIntegers.dvd_four_sub_sq_iff_even_even_of_ne_one_mod_four}
  \leanok
  If $d \not\equiv 1 \pmod 4$, then
  \[
  4 \mid (a'^2 - db'^2)
  \]
  if and only if both $a'$ and $b'$ are even.
\end{lemma}
\begin{proof}
  This is the key mod-$4$ arithmetic lemma for the non-\(1 \bmod 4\) branch.
\end{proof}

\begin{lemma}\label{bp_mod4_one}
  \lean{QuadraticNumberFields.RingOfIntegers.dvd_four_sub_sq_iff_same_parity_of_one_mod_four}
  \leanok
  If $d \equiv 1 \pmod 4$, then
  \[
  4 \mid (a'^2 - db'^2)
  \]
  if and only if $a'$ and $b'$ have the same parity.
\end{lemma}
\begin{proof}
  This is the key mod-$4$ arithmetic lemma for the \(1 \bmod 4\) branch.
\end{proof}

\begin{lemma}\label{bp_lift_nonone}
  \lean{QuadraticNumberFields.RingOfIntegers.exists_zsqrtd_of_isIntegral_of_ne_one_mod_four}
  \leanok
  \uses{bp_halfint, bp_mod4_nonone, bp_zsqrtd}
  If $d \not\equiv 1 \pmod 4$, every integral element of $\Q(\sqrt{d})$ comes
  from the model $\Z[\sqrt{d}]$.
\end{lemma}
\begin{proof}
  The half-integer normal form plus the non-\(1 \bmod 4\) parity lemma force
  both coordinates to be even, so the element lifts through
  \verb|Zsqrtd.toQsqrtdHom|.
\end{proof}

\begin{lemma}\label{bp_lift_one}
  \lean{QuadraticNumberFields.RingOfIntegers.exists_zOnePlusSqrtOverTwo_of_isIntegral_of_one_mod_four}
  \leanok
  \uses{bp_halfint, bp_mod4_one, bp_zone}
  If $d = 1 + 4k$, every integral element of $\Q(\sqrt{d})$ comes from the
  model $\Z\left[\frac{1+\sqrt{d}}{2}\right]$.
\end{lemma}
\begin{proof}
  The same-parity criterion is exactly the carrier condition for
  \verb|ZOnePlusSqrtOverTwo|.
\end{proof}

\begin{lemma}\label{bp_equiv_embedding}
  \lean{QuadraticNumberFields.RingOfIntegers.ringOfIntegers_equiv_of_embedding}
  \leanok
  A commutative ring embedding into a number field that captures precisely the
  integral elements is canonically isomorphic to the ring of integers.
\end{lemma}
\begin{proof}
  This packages the generic \verb|IsIntegralClosure| argument used by both branches.
\end{proof}

\begin{theorem}\label{bp_branch_nonone}
  \lean{QuadraticNumberFields.RingOfIntegers.ringOfIntegers_equiv_zsqrtd_of_mod_four_ne_one}
  \leanok
  \uses{bp_lift_nonone, bp_equiv_embedding}
  If $d \not\equiv 1 \pmod 4$, then
  \[
  \OO_{\Q(\sqrt{d})} \cong \Z[\sqrt{d}].
  \]
\end{theorem}
\begin{proof}
  The generic embedding criterion is applied to the canonical map from
  \verb|Zsqrtd d| into \verb|Qsqrtd (d : \Q)|.
\end{proof}

\begin{theorem}\label{bp_branch_one}
  \lean{QuadraticNumberFields.RingOfIntegers.ringOfIntegers_equiv_zOnePlusSqrtOverTwo_of_mod_four_eq_one}
  \leanok
  \uses{bp_lift_one, bp_equiv_embedding}
  If $d = 1 + 4k$, then
  \[
  \OO_{\Q(\sqrt{d})}
  \cong
  \Z\left[\frac{1+\sqrt{d}}{2}\right].
  \]
\end{theorem}
\begin{proof}
  The same generic embedding criterion is applied to the canonical map from
  \verb|ZOnePlusSqrtOverTwo k| into the field model.
\end{proof}

\begin{theorem}\label{bp_classification}
  \lean{QuadraticNumberFields.RingOfIntegers.ringOfIntegers_classification}
  \leanok
  \uses{bp_branch_nonone, bp_branch_one}
  The project proves the complete ring-of-integers classification by splitting
  on the two cases for $d \bmod 4$.
\end{theorem}
\begin{proof}
  This is the final disjunction theorem packaging both branch equivalences.
\end{proof}
